\documentclass[10pt,a4paper]{amsart}
\usepackage[margin=19mm]{geometry}
\usepackage{amsmath,amssymb,mathtools}
\usepackage[scr=rsfs,cal=euler]{mathalfa}
\usepackage{xfrac}
\usepackage[unicode,hidelinks,bookmarks=false]{hyperref}
\newcommand{\ie}{i.e.\ }
\newcommand{\cf}{cf.\ }
\newcommand{\resp}{respectively\ }
\newcommand{\Subsection}{\subsection}
\providecommand{\nobreakdash}{\nobreak\hbox{-}}
\setlength{\emergencystretch}{2em}
\multlinegap0pt
\usepackage{bm}
\usepackage{bbm}
\usepackage{dsfont}
\usepackage{xspace}
\usepackage{cancel}
\usepackage{mathtools}
\usepackage{amsthm}
\makeatletter
\@ifundefined{assumption}{\newtheorem{assumption}{Assumption}}{}
\@ifundefined{theorem}{\newtheorem{theorem}{Theorem}}{}
\renewenvironment{proof}[1]{\par\medskip\noindent\textbf{#1}\enspace\ignorespaces}{\par\medskip}
\makeatother
\title[Model-Driven Subspaces for Large-Scale Optimization with Loc]{Model-Driven Subspaces for Large-Scale Optimization with Local Approximation Strategy}
\author{Yitong He, Pengcheng Xie}
\date{Source: arXiv:2509.08256v1 (CC BY 4.0)}
\begin{document}
\maketitle

\noindent\emph{Source and licence.} Model-Driven Subspaces for Large-Scale Optimization with Local Approximation Strategy. Version arXiv:2509.08256v1 (CC BY 4.0), distributed under the Creative Commons Attribution 4.0 International licence, as recorded in the arXiv licence metadata for this exact version and shown on the abstract page https://arxiv.org/abs/2509.08256v1. Full rights notice: licenses/arxiv-2509.08256-CC-BY-4.0.txt.\par\smallskip

\noindent\textit{Editorial scope.} Counted unit: the Theorem whose source label is \texttt{ineq:pre\_cauchy decrease} in the article (source0.tex lines 666--672, source label \texttt{ineq:pre\_cauchy decrease}), delivered together with the article's own complete original proof of it (source0.tex lines 674--734, one \texttt{proof} environment of 61 lines). No mathematics is written, repaired or re-derived: statement and proof are the article's own text line for line. Required context retained uncounted: ineq:1/2 decrease of our algo (lines 634--636); assumption: cauchy (lines 655--662). Every definition, display and label of those lines is retained unchanged. Omitted from this package and not redistributed: 1--633, 637--654, 663--665, 673--673, 735--1263. Nothing inside a retained excerpt is omitted or altered: every retained body is delivered exactly as the article's lines are written, so no omission receipt of any kind accompanies this package. Citations: the retained excerpts carry no citation command at all, so no bibliography entry of the article is transcribed and no reference is redistributed. Reference adaptation: none was needed and none was made; every reference inside the delivered unit resolves inside it, so no unresolved reference or citation is produced. Printed slips kept literal: no printed word, symbol, hypothesis or number of the counted statement or of its proof was altered. Layout and class: the article's own class is \texttt{article} with packages including orcidlink, booktabs, doi, geometry, boxedminipage, epsfig, amssymb, amsmath, version, graphicx, fancybox, mathrsfs, bm, amsthm; the wrapper is the lane's pinned \texttt{amsart} editorial wrapper at 10pt on A4, which restates the theorem-like environments the retained text needs and adds the article's own macro definitions that the retained text needs (none), with extra wrapper packages (bm, bbm, dsfont, xspace, cancel, mathtools). The delivered numbering is the unit's own. Assets: no retained span contains a figure environment, a graphics-inclusion command, a PDF-page inclusion, a table or a TikZ picture; no image file is copied into this package, the package declares no asset dependency and no image file is redistributed with it. Licence: version arXiv:2509.08256v1 of this article is distributed under the Creative Commons Attribution 4.0 International licence, as recorded in the arXiv licence metadata for this exact version and shown on the abstract page https://arxiv.org/abs/2509.08256v1. A full rights notice is delivered at licenses/arxiv-2509.08256-CC-BY-4.0.txt. Characters: every retained excerpt is pure ASCII, so the delivered engine, XeLaTeX with its default Latin Modern text font, reports no missing character.
\begin{equation}\label{ineq:1/2 decrease of our algo}
    m_k(\bm 0)-m_k(\bm y^*)\geq \frac{1}{2}\left(m_k(\bm 0)-m_k(\bm s^*)\right).
\end{equation}

\begin{assumption}\label{assumption: cauchy}
    The Hessian $\nabla^2m_k(\bm 0)$ of model $m_k(\bm s)$
    is positive definite. Denote $q=\dim(\mathcal{S}_k)$.
    Let $\lambda_1,\cdots,\lambda_q$ be the eigenvalues of $\nabla^2m_k(\bm 0)$, then 
    \begin{equation*}  \left(\sum_{i=1}^q\lambda_i^2\right)\left(\sum_{i=1}^q\frac{1}{\lambda_i}\right)\leq c,
    \end{equation*}
    where $c\in[1/4,1/2]$ is a constant independent of $k$.
\end{assumption}

\begin{theorem}
Suppose $m_k$ satisfies Assumption \ref{assumption: cauchy} and \eqref{ineq:1/2 decrease of our algo}, then the computed solution $\bm y^*$ satisfies
\begin{equation}\label{ineq:pre_cauchy decrease}
    m_k(\bm 0)-m_k(\bm y^*)\geq c_1\Vert\nabla m_k(\bm 0)\Vert^2
\end{equation}
for some constant $c_1\in[1/2,1]$ independent of $k$.
\end{theorem}

\begin{proof}{Proof}
    Since $m_k$ satisfies \eqref{ineq:1/2 decrease of our algo}, \begin{equation*}
    \begin{aligned}
        m_k(\bm 0)-m_k(\bm y^*)
        &\geq\frac{1}{2}\left((m_k(\bm 0)-m_k(\bm s^*)\right)\\
        &=-\frac{1}{2}\left(\bm g^\top \bm s^*+\frac{1}{2}\bm s^{*\top}\bm B\bm s^*\right)\\
        &=-\frac{1}{2}\left(\nabla m_k(\bm s^*)-\frac{1}{2}\bm B\bm s^*\right)^\top\bm s^*,
    \end{aligned}
    \end{equation*}
where $\bm s^*$ is the global minimizer of $m_k$. Therefore $\nabla m_k(\bm s^*)=\bm 0$ and 
\begin{equation*}
        -\frac{1}{2}\left(\nabla m_k(\bm s^*)-\frac{1}{2}\bm B_k\bm s^*\right)^\top\bm s^*=\frac{1}{4}\bm s^{*\top}\bm B_k\bm s^*.
\end{equation*}
By our assumption, $\bm B_k=\nabla^2m_k(\bm 0)=\nabla^2m_k(\bm s^*)$ is positive definite. Hence there is a orthogonal matrix $\bm Q\in\mathbb{R}^{q\times q}$ such that
\begin{equation*}
    \bm Q^\top \bm B_k\bm Q={\rm diag}\{\lambda_1,\cdots,\lambda_q\},\ \lambda_i>0.
\end{equation*}

Let $\bm s^*=\bm Q \bm s$, $\bm g_k=\bm Q\bm y$, where $\bm s=(s_1,\cdots,s_q)^\top$ and $\bm y=(y_1,\cdots,y_q)^\top$, then 
\begin{equation*}
\begin{aligned}
     \frac{1}{4}\bm s^{*\top}\bm B_k\bm s^*=\frac{1}{4}\bm s^\top\bm Q^\top\bm B_k\bm Q\bm s=\frac{1}{4}\sum_{i=1}^{q}\lambda_i s_i^2.
\end{aligned}
\end{equation*}
By Cauchy's inequality, we have 
\begin{equation*}
    \sum_{i=1}^{q}\lambda_i s_i^2\geq \frac{\left(\sum_i^q| s_i|\right)^2}{\sum_{i=1}^q\frac{1}{\lambda_i}}\geq\frac{\Vert\bm s\Vert^2}{\sum_{i=1}^q\frac{1}{\lambda_i}}.
\end{equation*}
Hence 
\begin{equation*}
\begin{aligned}
     \frac{1}{4}\sum_{i=1}^{q}\lambda_i s_i^2\geq\frac{1}{4}\frac{\Vert\bm s\Vert^2}{\sum_{i=1}^q\frac{1}{\lambda_i}}=\frac{1}{4\sum_{i=1}^q\frac{1}{\lambda_i}}\Vert\bm Q^\top \bm s^*\Vert.
\end{aligned}
\end{equation*} 
Moreover, since $\Vert \bm Q^\top\bm s^*\Vert=\Vert \bm Q^\top\bm B_k^{-1}(-\bm g_k)\Vert=\sum_{i=1}^q y_i^2/\lambda_i^2$, we have
\begin{equation*}
    \frac{1}{4}\sum_{i=1}^{q}\lambda_i s_i^2\geq\frac{1}{4\sum_{i=1}^q\frac{1}{\lambda_i}}\Vert\bm Q^\top \bm B_k^{-1}\left(-\bm g_k\right)\Vert=\frac{1}{4\sum_{i=1}^q\frac{1}{\lambda_i}}\sum_{i=1}^q \left(\frac{y_i}{\lambda_i}\right)^2.
\end{equation*}

Again, applying Cauchy's inequality, we have
\begin{equation*}
    \sum_{i=1}^q \left(\frac{y_i}{\lambda_i}\right)^2\geq \frac{(\sum_i^q|y_i|)^2}{\sum_{i=1}^q\lambda_i^2}\geq \frac{\Vert\bm y\Vert^2}{\sum_{i=1}^q\lambda_i^2}.
\end{equation*} 
Therefore
\begin{equation*}
\begin{aligned}
    \frac{1}{4\sum_{i=1}^q\frac{1}{\lambda_i}}\sum_{i=1}^q \left(\frac{y_i}{\lambda_i}\right)^2\geq \frac{1}{4\sum_{i=1}^q\frac{1}{\lambda_i}}\frac{\Vert\bm y\Vert^2}{\sum_{i=1}^q\lambda_i^2}\geq \frac{1}{4\sum_{i=1}^q\frac{1}{\lambda_i}}\frac{\Vert\bm g_k\Vert^2}{\sum_{i=1}^q\lambda_i^2},
\end{aligned}  
\end{equation*}
where the second inequality comes from 
\begin{equation*}
    \Vert \bm g_k\Vert=\Vert\bm Q\bm y\Vert\leq \Vert\bm Q\Vert\Vert\bm y\Vert=\Vert\bm y\Vert,
\end{equation*}
$\Vert\bm Q\Vert=\sup_{\Vert\bm x\Vert=1}\Vert\bm Q\bm x\Vert$ is the matrix norm of $\bm Q$ and $\Vert\bm Q\Vert=1$ since $\bm Q$ is an orthogonal matrix. 

Let $c_1=\frac{1}{4c}$, then by Assumption \ref{assumption: cauchy}, $c_1\in[1/2,1]$ and 
\begin{equation*}
     c_1\leq\frac{1}{4}\frac{1}{\sum_{i=1}^m\frac{1}{\lambda_i}\sum_{i=1}^m\lambda_i^2}.
\end{equation*}
The proof is completed by observing $\bm g_k=\nabla m_k(\bm 0)$.  
\end{proof}

\end{document}
